\VerdierPhysicalPageAnchor{0019}%
gramme ci-après soit commutatif :

\VerdierSemanticAnchor{intro-1-resolution-comparison-diagram}
\[
\begin{tikzcd}[column sep=large,row sep=large]
X \arrow[r,"\varepsilon'"] \arrow[dr,"\varepsilon"']
  & \mathcal R(X) \arrow[d,"\rho"] \\
  & I^{\bullet}(X)
\end{tikzcd}
\]

\VerdierSemanticAnchor{intro-1-f-acyclic-morphism}
De plus, étant donné un morphisme de résolutions $F$-acycliques :

\VerdierSemanticAnchor{intro-1-f-acyclic-morphism-diagram}
\[
\begin{tikzcd}[column sep=large,row sep=large]
X \arrow[r,"\varepsilon_1"] \arrow[dr,"\varepsilon_2"']
  & \mathcal R_1(X) \arrow[d,"\rho'"] \\
  & \mathcal R_2(X)
\end{tikzcd}
\qquad\text{,}
\]
le morphisme de complexes d'objets de $\mathcal B$ :

\VerdierSemanticAnchor{intro-1-f-rho-prime}
\[
F(\rho') : F(\mathcal R_1(X)) \longrightarrow F(\mathcal R_2(X))
\]
est un quasi-isomorphisme.

\VerdierSemanticAnchor{intro-1-inductive-system}
Remarquons qu'ici on obtient dans $\mathcal B$ un système inductif de
complexes et de quasi-isomorphismes qui possède un élément maximal unique à
homotopie près : le transformé par $F$ d'une résolution injective de $X$. Ce
phénomène ne se produit plus dans le deuxième exemple que nous abordons
maintenant.

\VerdierSemanticAnchor{intro-1-left-derived-tensor}
Le procédé de définition des foncteurs dérivés par résolution se généralise
parfois au cas où la catégorie $\mathcal A$ ne possède plus nécessairement
suffisamment d'objets injectifs ou projectifs. Considérons le cas où
$\mathcal A$ est la catégorie des faisceaux de groupes commutatifs sur un
espace topologique $E$. Cette catégorie ne possède pas, en général,
suffisamment d'objets projectifs. On se propose néanmoins de définir les
\emph{foncteurs dérivés gauches} du foncteur exact à droite de $\mathcal A$
dans $\mathcal A$ qui, à tout faisceau $G$, associe le faisceau
$F\otimes_{\mathbb Z}G$ ($F$ faisceau fixe). On exige, bien entendu, que cette
définition fournisse des foncteurs possédant les propriétés usuelles des
foncteurs dérivés gauches (comportement par rapport aux suites exactes et une
certaine propriété universelle) et qu'elle redonne dans les cas connus (par
exemple lorsque $E$ est réduit à un point) les produits de torsion usuels.

\VerdierSemanticAnchor{intro-1-left-resolution}
On peut opérer de la manière que voici. Tout faisceau $G$ admet une résolution
à gauche $L^{\bullet}\to G$ de longueur finie (en fait de longueur 2) par des
